\chapter{Mathematical Foundations}
\label{ch:essay-math}

The four theorems developed in \emph{Admissible Reconstruction} depend on a common apparatus that can be stated concisely. An admissibility structure $R$ on a set of states $X$ determines, for any state $x$, a reachable set $\reach{x}$ consisting of all states $y$ such that $x \adm{R} y$. Nothing about $R$ is assumed beyond this: it is a relation, not a metric, and it need not be symmetric, transitive, or total. The relation may vary with time, written $R_t$, so that what counts as admissible continuation changes as the system evolves. The four theorems below are each derivable from this minimal starting point, with no additional primitive vocabulary required.

\section{Representation Invariance}

Two representations built from different substrates are semantically equivalent when their admissibility structures are related by an admissibility-isomorphism. A map $\varphi : X \to X'$ between state spaces $(X, R)$ and $(X', R')$ is an admissibility-isomorphism when, for all states $x, y \in X$, the biconditional $x \adm{R} y$ if and only if $\varphi(x) \adm{R'} \varphi(y)$ holds. The definition is deliberately minimal: $\varphi$ preserves the admissibility relation in both directions, and nothing more is required.

Theorem~1.1 of \emph{Admissible Reconstruction}, Chapter~1, states that if such a $\varphi$ exists, then $(X, R)$ and $(X', R')$ are semantically equivalent, regardless of the physical or computational substrate realizing either representation. The theorem is close to definitional, and this is by design. Semantic equivalence within the admissibility framework has never been defined as identity of physical encoding but as identity of what a representation permits and forbids as continuation. The theorem makes this definition explicit as an invariance statement rather than leaving it distributed across the apparatus.

What is not definitional is establishing that a particular map $\varphi$ between two representations of interest actually satisfies the biconditional in both directions. This is the harder direction in practice: showing that a transformation preserves reachable-set structure exactly, rather than merely preserving some coarser statistic that happens to look similar. A claim that speech and gesture are representationally invariant, for instance, requires constructing an explicit $\varphi$ and verifying that admissible continuations in one modality correspond exactly to admissible continuations in the other, with neither modality admitting a continuation the other forbids.

Corollary~1.2 of Chapter~1, which is essential to the cross-modal argument developed in later sections, establishes that no property of the state space that is not preserved by every admissibility-isomorphism can be a semantic property of the representation. In particular, substrate-specific properties such as the physical channel through which a representation is realized, its production cost, or its latency are not semantic properties unless they happen to be expressible as constraints on the admissibility relation itself. This does not render substrate properties unimportant; it provides a criterion for deciding when they are semantically relevant. If changing a physical property leaves the admissibility structure invariant, that property does not distinguish the representations at the level recognized by the theory. If changing it alters the reachable relations, then it has become structurally relevant. The corollary thereby draws a sharp boundary between what survives representational equivalence and what belongs to the study of implementation rather than meaning.

\section{Admissibility Preservation}

A trajectory is a sequence of states $S_0, S_1, \dots, S_t$ such that $S_i \adm{R} S_{i+1}$ for each step, where $R$ may itself vary with time. The trajectory degenerates at step $i$ if the reachable set $\reach{S_i}$ is either empty, so that no continuation is admissible, or a singleton entirely determined by the history $S_0, \dots, S_{i-1}$, so that exactly one continuation is admissible and it was already fixed before $S_i$ was reached. In the first degenerate case, the admissibility relation restricted to continuations from $S_i$ is uniformly false; in the second, it is uniformly true of exactly one candidate and uniformly false of every other. Either way, the relation is doing no discriminating work at step $i$.

Theorem~2.1 of Chapter~2 states that a trajectory is coherent if and only if it does not degenerate at any step. Equivalently, the trajectory is coherent if and only if the reachable-set cardinality $|\reach{S_i}|$ exceeds a nontrivial floor relative to what has already been committed at every step along the way. Coherence has sometimes been treated informally as though it were an additional property a trajectory could possess independent of the admissibility relation itself, as though a trajectory could be admissible at every step and yet somehow incoherent, or coherent despite violating admissibility. The theorem rules this out: coherence and nondegeneracy of the admissibility relation are the same property under two names.

Corollary~2.1, proven in the same chapter, addresses operational tests for coherence. Domain-specific measures such as minimum support ratios required of a discourse trajectory, consistency scores, and similar quantities are frequently introduced as though they were primitive definitions of coherence in their own right. The corollary establishes that any such test is correctly understood only as an operational proxy for the condition that $|\reach{S_i}|$ remains above a nontrivial floor, not as an independent criterion. Where a proposed proxy can be false while the reachable set remains nontrivial, or true while the reachable set has degenerated, the proxy has diverged from what coherence actually consists in and should be revised accordingly.

\section{Monotonic Relaxation}

A scalar functional $\closure{X} = -\log \mu(\reach{X})$, where $\mu$ is a measure on the state space assigning nonzero mass to nonempty reachable sets, provides a measure of inadmissibility: the closure functional is large when few futures remain reachable from $X$ and small when many remain. A trajectory follows repair dynamics when each transition $S_i \to S_{i+1}$ is selected, among admissible continuations, to increase reachable-set volume: $\mu(\reach{S_{i+1}}) \ge \mu(\reach{S_i})$. This characterization is inherited from the restorative branch of the modificatory-operation taxonomy developed elsewhere in this research program, where repair is already defined as the operation type that restores lost admissible continuations.

Theorem~3.1 of Chapter~3 states that if a trajectory follows repair dynamics throughout, then the closure functional is nonincreasing along the trajectory: $\closure{S_{i+1}} \le \closure{S_i}$ for every step. The proof is immediate once repair is defined as volume-increasing, since $-\log$ is a strictly decreasing function on positive reals. The theorem's content is therefore not in the derivation but in the definition of repair dynamics being the correct operational characterization. The hypothesis is essential. Constitutive and regulatory modifications, the other two branches of the operation taxonomy, need not increase reachable-set volume, and ordinary state evolution under a fixed admissibility structure can move toward states with strictly smaller reachable sets without this constituting any failure of admissibility. A claim of the form ``inadmissibility always decreases'' would be false. Theorem~3.1 holds only of the restorative subclass, and its scope must be stated with this restriction.

The theorem shares a name with monotonic continuation, an ordering relation $S_t \le S_{t+1}$ developed elsewhere in this program, but the two claims are distinct and should not be conflated. Monotonic continuation is an ordering asserting that later states preserve the distinctions from earlier states that remain necessary for the reachable trajectory; it says nothing about any scalar quantity increasing or decreasing. Monotonic relaxation is a claim that a specific scalar functional decreases along a specific class of trajectories. Corollary~3.1 of Chapter~3 establishes the precise relationship: the two claims coincide only under an additional hypothesis that must be checked independently, and neither implies the other in general. A trajectory can increase reachable-set volume while failing to retain the specific significant continuations required by monotonic continuation, by replacing them with an equal or greater volume of differently located reachable states. Conversely, a trajectory can satisfy monotonic continuation's retention condition while reachable-set volume decreases, if what is retained is a shrinking but still significant core. The corollary exists to prevent the two theorems from being cited interchangeably.

\section{Non-Identifiability}

An observation map $\obsmap : \History \to \Observable$ takes complete histories to whatever observable evidence remains available to an observer without access to the history itself. The elementary direction of the non-identifiability result is nearly definitional. If two distinct histories $H_1 \neq H_2$ satisfy $\obsmap(H_1) = \obsmap(H_2)$, then no function $f : \Observable \to \History$ can satisfy both $f(\obsmap(H_1)) = H_1$ and $f(\obsmap(H_2)) = H_2$ simultaneously, since $f$ must be single-valued. This establishes impossibility once a collision has already been shown to exist, but it says nothing about whether such collisions are rare pathological cases or generic features of realistic observation maps.

The structural direction, Theorem~4.2 of Chapter~4, addresses this directly. Consider any observation map that factors through the reachable-set structure, meaning that $\obsmap(H) = \psi(\reach{H})$ for some function $\psi$, so that the observation map records only which continuations remain admissible from $H$ rather than the full state of $H$ itself. Whenever two histories $H_1 \neq H_2$ satisfy $\reach{H_1} = \reach{H_2}$, it follows immediately that $\obsmap(H_1) = \obsmap(H_2)$, and the elementary direction applies. The theorem's substantive claim is that this is not a degenerate case. The reachable-set operator $R(\cdot)$ partitions the space of histories into equivalence classes by definition: two histories lie in the same class exactly when they share a reachable set. An admissibility structure has a trivial quotient only when no two distinct histories ever share a reachable set, which would require the partition to be discrete. This is not the generic case for admissibility structures where reachable sets are determined by bounded information about the history while the history itself can vary over an unbounded space. The pigeonhole principle then guarantees that histories differing in unrecorded portions while agreeing in recorded portions will share a reachable set and thereby collide under the observation map. Non-injectivity is therefore a structural consequence of the observation map's design, not a contingent misfortune.

Corollary~4.1, the diagnostic corollary, establishes that the question of why reconstruction fails in a given case admits a principled decomposition. Given an observed collision $\obsmap(H_1) = \obsmap(H_2)$ with $H_1 \neq H_2$, one can check whether $\reach{H_1} = \reach{H_2}$. If the reachable sets differ, then the observation map has discarded information that a less lossy procedure could in principle have retained, and the failure is a correctable error of the particular observation procedure. If the reachable sets are equal, then no observation map factoring through reachable-set structure can distinguish $H_1$ from $H_2$, and the failure is a structural limit rather than a procedural defect. The corollary thereby provides a criterion for deciding whether a specific case of reconstruction uncertainty is correctable or fundamental, and this criterion is in principle decidable by examining the reachable-set structure directly.

Together, these four theorems establish that representation, coherence, repair, and the limits of recovery are not four independent subjects requiring separate theoretical machinery. They are four aspects of a single admissibility relation, examined at different points in its operation: Representation Invariance at a single state, Admissibility Preservation along a trajectory, Monotonic Relaxation under restorative dynamics, and Non-Identifiability under lossy projection.
